\chapter*{Avant-propos}
\addcontentsline{toc}{chapter}{Avant-propos}

Ce rapport présente mon travail de Bachelor réalisé à la HEIG-VD : la conception d'un lecteur audio portable, du schéma électronique au firmware.

J'ai choisi ce sujet parce qu'il réunit des domaines habituellement traités séparément au cours de la formation, à savoir l'électronique analogique et numérique, l'alimentation, le routage d'une carte multicouche et le logiciel embarqué. La difficulté principale s'est concentrée sur la gestion des interfaces entre ces parties.

Au-delà du cadre académique, ce projet est avant tout personnel. Mon objectif n'était pas de produire une étude ou une démonstration de principe, mais un appareil réel, fabriqué, assemblé et fonctionnel, que l'on peut tenir en main et utiliser. Cette exigence a guidé l'ensemble des décisions, du choix des composants à l'organisation du firmware, car un appareil destiné à durer ne tolère pas les approximations qu'une maquette admettrait.

Le document suppose des bases en électronique et en systèmes embarqués. Je me suis attaché à justifier les choix de conception plutôt qu'à seulement les énoncer.